\chapter{The Data page}\label{ch:data}

The \ui{Data} page is where recorded data are downloaded, inspected and exported. The workflow is always the same: \textbf{pick a
station, pick a time range, choose how to process it, press the big button, then look and export.}

\begin{figure}[tbp]
\centering
\annotshot[\linewidth]{data-top}{\calloutsData}
\caption{The Data page before a preview. (1) Station. (2) Import a miniSEED file. (3) Start and end. (4) Presets.
(5) UTC or local time. (6) Visible span of the coverage bar. (7) Move the visible span. (8) Coverage bar. (9) Selection handle.
(10) Files that will be downloaded. (11) Channels. (12) Band-pass. (13) Units. (14) Download and preview.}\label{fig:datatop}
\end{figure}

\section{Choose a station}
The \ui{Station} list (1) contains your Shake, and every other station whose data are on this laptop (for example from imported
files). A station name has three parts separated by dots: network, station and location, for example \code{AM.R77A7.00}.

\section{Choose a time range}
\paragraph{Typing.} The \ui{Start} and \ui{End} fields (3) take a date and time written as year-month-day hour:minute:second, for example \code{2026-10-15 03:00:00}. Press
\key{Enter} or click away to apply. A malformed value shows a hint with the expected form. The text next to the toggle shows
the length of the range.

\paragraph{Presets.} The clock button (4) opens a menu of ready-made ranges.

\begin{figure}[tbp]
\centering
\shot[0.8\linewidth]{data-presets}
\caption{The presets menu: \ui{Last 10 min}, \ui{Last hour}, \ui{Today}, \ui{Yesterday} and \ui{Last 7 days}.}
\end{figure}

\paragraph{UTC or local.} Seismic data are always stored and exported in UTC. The toggle (5) only changes the time zone in which the
fields and axes are \emph{displayed and typed}; \ui{Today} and \ui{Yesterday} follow it. The code and the files you export always use UTC.

\section{The coverage bar}
The coloured bar (8) shows what exists for the visible span.
\begin{center}\small
\begin{tabular}{@{}cl@{}}
\textcolor{rs3dgreen}{\rule{1.2em}{1ex}} & \textbf{green}: the data are on this laptop\\
\textcolor{rs3dgrey}{\rule{1.2em}{1ex}} & \textbf{grey}: the data exist only on the Shake and have not been downloaded\\
\textcolor{rs3dred}{\rule{1.2em}{1ex}} & \textbf{red}: no data anywhere (a gap, or time that has not happened yet appears blank)\\
\end{tabular}
\end{center}

The strip below it (9) is the selection. Drag in it to select a range, drag its handles to change the start or end, or drag the
middle to move it; the \ui{Start} and \ui{End} fields follow. \ui{Week}, \ui{Day} and \ui{Hour} (6) change how much time the bar shows, and
the arrows (7) move to the previous or next span.

\begin{figure}[tbp]
\centering
\shot[0.95\linewidth]{data-coverage}
\caption{The coverage bar for 15~October: green where the demo data are on the laptop, and the selection (light blue) about 03:00--04:00.}
\end{figure}

\paragraph{Files to download.} The number next to (10) is how many day files will be copied from the Shake when you press the
button. Today's file is always copied again, because the Shake keeps writing to it.

\section{Choose channels, filter and units}
\paragraph{Channels (11).} Tick the channels to load: \code{EHZ} (vertical), \code{EHN} (north--south), \code{EHE} (east--west).

\paragraph{Band-pass (12).} Switch it on to keep only the frequencies between the two boxes (1 and 10~Hz at first) and remove
drift and high-frequency noise. The band you enter is remembered for next time. The filter is zero-phase, so it does not delay the
waveform, but it smears sharp onsets slightly backwards in time; do not use it to read exact arrival times (Theory Manual, Chapter~4).

\paragraph{Units (13).}\label{sec:units}
\ui{counts} are the raw numbers of the converter. \ui{m/s} converts them to ground velocity using the instrument response. The first
time you choose \ui{m/s} for a station the application downloads the response once from the Raspberry Shake service (the laptop
needs Internet that one time) and keeps it, after which the switch works offline. Without it and without Internet you get an error
message instead of a chart.

\section{Download and preview}
Press the large button (14). For your own Shake it reads \ui{Download \& preview}: it first copies the day files you do not yet
have, showing progress (\emph{Downloading 2/6: EHZ 2026-10-14}), then draws the data. For a station that is not a connected Shake
(for example an imported file) the button reads \ui{Preview} and only reads local files. If the Shake is switched off the
download is skipped, an explanation is shown and whatever is already on the laptop is displayed.

\section{Reading the preview}
\begin{figure}[tbp]
\centering
\shot[\linewidth]{data-waveforms}
\caption{The waveform preview of one hour (03:00--04:00): three channels sharing one time axis. The short spike near 03:12 is the
local earthquake; the amplitude axes are in raw counts.}\label{fig:datawf}
\end{figure}

\paragraph{Waveforms.} One plot per selected channel, with the time axis in the zone you chose. Hover for the value at a point. The
counts do not oscillate around zero because they sit on the instrument's offset (Theory Manual, Section~3.4); switching the
band-pass on removes it.

\paragraph{Data health.} Just under the export buttons, one line says how complete the selected range is and how many gaps it has,
for example \emph{100.00\,\% complete $\cdot$ no gaps}. If there are gaps their start, end and length are listed in the time
zone you chose (Theory Manual, Chapter~8).

\paragraph{Spectrogram.} A picture of \emph{which frequencies were present when}. Time runs left to right, frequency (0 to 50~Hz)
bottom to top, colour is power in decibels (red strong, blue weak). Use the channel selector next to the title to choose
\code{EHZ}, \code{EHN} or \code{EHE}. An earthquake is a vertical streak; a steady horizontal line is a constant interfering
frequency, here a 30~Hz line that is part of the demonstration data. Blank columns are gaps.

\begin{figure}[tbp]
\centering
\shot[\linewidth]{data-spectrogram}
\caption{The spectrogram of the same hour for \code{EHZ}. The horizontal red line near 30~Hz is a steady interference; the dark
patch near 03:12 is the earthquake.}
\end{figure}

\paragraph{Noise spectrum.} The median level of each frequency over the whole selection, with a band showing the 10th to 90th
percentile. It is the quickest way to compare a quiet period with a busy one: select night, preview, select midday, preview again.
It is in raw units and is meant for \emph{comparing your own selections}, not for comparing with published noise models
(Theory Manual, Chapter~7).

\begin{figure}[tbp]
\centering
\shot[\linewidth]{data-noise}
\caption{The noise spectrum: median (green) and 10th/90th percentile (dashed). The spike at 30~Hz is the interference line; the fall
near 50~Hz is the Shake's anti-alias filter.}
\end{figure}

\section{Exporting}
After a preview the export buttons below the waveforms become active: \ui{miniSEED}, \ui{CSV}, \ui{PNG}, \ui{PDF} and
\ui{</> Code}. The file name is the network, the station, the start and the end time in UTC, and the extension. For example:
\begin{center}\code{AM.DEMO1\_20261015T030000\_20261015T040000.mseed}\end{center}
Exports contain exactly what the preview shows: if the
band-pass was on, the file is filtered; if you chose \ui{m/s}, it is in m/s.

\begin{description}[leftmargin=3.2cm,style=nextline]
\item[miniSEED] the standard seismological format, readable by ObsPy, SeisComP and most other tools. Gaps are kept.
\item[CSV] a text table: one column of UTC times and one column per channel, one row per sample; missing samples are empty.
  All channels must have the same sampling rate. See Appendix~\ref{app:formats}.
\item[PNG, PDF] a picture of the waveforms for reports.
\item[</> Code] opens a dialog with a short Python (ObsPy) program that reproduces \emph{exactly} the processing behind the preview,
  and a button to save it as a \file{.py} file. It reads from this laptop's data folder, so it runs without the Shake. Use it when
  you want to go further than the application does, or to document how a number was obtained.
\end{description}

\begin{figure}[tbp]
\centering
\shot[0.82\linewidth]{data-code}
\caption{The generated Python program for the preview above.}
\end{figure}

\section{Importing miniSEED files}
\ui{Import miniSEED} (2) adds a file from anywhere, for example from another Shake, a university network or a data centre. Accepted
names end in \code{.mseed}, \code{.msd}, \code{.miniseed} or \code{.ms}, up to 2~GB. The file is split into one file per channel
per day in your data folder, next to the data from your Shake, and the page switches to the imported station. Importing the same
file twice changes nothing; importing a file with overlapping data merges them. Imported stations then appear in the \ui{Station}
list, on the \ui{Helicorder} page and everywhere else.

\begin{rs3dtip}[No Shake needed]
Import is also the quickest way to try the application: the whole Data and Helicorder pages work on any miniSEED file.
\end{rs3dtip}
